\sectionkicker{Source-backed technical notes}
\subsection*{汎用Agent競争の外側 — Deep Thinkと科学推論: Technical Notes}
本欄は記事本文で圧縮した一次資料上の情報を、比較・再検証しやすい形で整理したものである。時系列、確認済みの主張、留意点、一次資料URLを掲載する。
\medskip
\subsection*{Theme at a glance}
\addcontentsline{toc}{subsection}{Theme at a glance}
\begin{center}
\footnotesize
\begin{tabularx}{\linewidth}{@{}p{0.20\linewidth}p{0.10\linewidth}p{0.14\linewidth}X@{}}
\toprule
資料 & 位置づけ & 種別 & 時系列 \\
\midrule
Gemini 3 Deep Think — February 2026 update & 主要資料 & モデル更新 & 2026-02-12 (モデル更新) \\
Gemini 3.1 Pro / Gemini 3.1 Flash Image February API releases & 補足資料 & モデル更新 & 2026-02-19 (API\_モデル公開); 2026-02-26 (API\_モデル公開) \\
\bottomrule
\end{tabularx}
\end{center}
\normalsize

\begin{technicalnote}{Gemini 3 Deep Think — February 2026 update}{主要資料}
\begingroup
% reader-facing Technical Notes break policy
\widowpenalty=10000
\clubpenalty=10000
\displaywidowpenalty=10000
\begin{tabularx}{\linewidth}{@{}>{\bfseries}p{0.22\linewidth}X@{}}
組織 & Google DeepMind \\
種別 & モデル更新 \\
時系列 & 2026-02-12 (モデル更新) \\
\end{tabularx}
\smallskip
{\bfseries 一次資料から整理したtechnical points}
\begin{itemize}[leftmargin=1.5em,itemsep=0.35em]
\item \textbf{一次情報で確認できる事実}: Google DeepMindは2026年2月12日、science、research、engineeringの問題を対象とするspecialized reasoning modeとして、更新版Gemini 3 Deep Thinkを発表した。
\item \textbf{一次情報で確認できる事実}: 2月の発表では、Google AI Ultra subscriber向けのGemini app提供とAPI利用のearly-access pathが説明され、広範な一般API availabilityではなかった。
\item \textbf{Vendor claim}: 発表に含まれる科学領域およびbenchmarkの性能結果は、独立再現されない限りvendor claimとして扱う。
\end{itemize}
\begin{minipage}{\linewidth}
% reader-facing Technical Notes generic boundary/source tail group
{\bfseries 読む際の境界}
\begin{itemize}[leftmargin=1.5em,itemsep=0.35em]
\item \textbf{分析上の留意点}: Deep Thinkはspecialized reasoning modeであり、その結果をGemini 3 family全体へ一般化すべきではない。
\item \textbf{分析上の留意点}: 発表時点のAPI availabilityは限定的なearly accessだった。
\item \textbf{分析上の留意点}: 科学benchmarkおよび比較に関する主張はvendorによる報告である。
\end{itemize}
\begin{samepage}
{\bfseries 一次資料}
\begingroup\sloppy
\Urlmuskip=0mu plus 2mu\relax
\begin{itemize}[leftmargin=1.5em,itemsep=0.25em]
\item {\scriptsize\url{https://deepmind.google/blog/}}
\end{itemize}
\endgroup
\end{samepage}
\end{minipage}
\endgroup
\end{technicalnote}
\medskip

\begin{technicalnote}{Gemini 3.1 Pro / Gemini 3.1 Flash Image February API releases}{補足資料}
\begingroup
% reader-facing Technical Notes break policy
\widowpenalty=10000
\clubpenalty=10000
\displaywidowpenalty=10000
\begin{tabularx}{\linewidth}{@{}>{\bfseries}p{0.22\linewidth}X@{}}
組織 & Google \\
種別 & モデル更新 \\
時系列 & 2026-02-19 (API\_モデル公開); 2026-02-26 (API\_モデル公開) \\
\end{tabularx}
\smallskip
{\bfseries 一次資料から整理したtechnical points}
\begin{itemize}[leftmargin=1.5em,itemsep=0.35em]
\item \textbf{一次情報で確認できる事実}: Gemini API changelogには、2026年2月19日のGemini 3.1 Pro Previewと、built-in execution capabilityと併せてuser-defined toolを優先できるcustom-tools endpointが記録されている。
\item \textbf{一次情報で確認できる事実}: changelogには、Gemini 3.1 Flash Image Preview（Nano Banana 2）が2026年2月26日に、高速・高volumeのimage generation向けreleaseとして記録されている。
\item \textbf{Vendor claim}: これらのpreviewに関連する性能・品質の優位性主張はvendor claimとして扱う。
\end{itemize}
\begin{minipage}{\linewidth}
% reader-facing Technical Notes generic boundary/source tail group
{\bfseries 読む際の境界}
\begin{itemize}[leftmargin=1.5em,itemsep=0.35em]
\item \textbf{分析上の留意点}: 両model identifierはいずれも2月時点ではpreview-stageのreleaseであり、その時点で一般にstableだったとは記述しない。
\item \textbf{分析上の留意点}: API changelogはchronologyのEvidenceであり、品質やbenchmark claimを独立に検証するものではない。
\item \textbf{分析上の留意点}: このtaskはreasoning/tool-use modelの更新とimage-generation modelの更新を含むため、編集上は両者の役割を必要に応じて分離する。
\end{itemize}
\begin{samepage}
{\bfseries 一次資料}
\begingroup\sloppy
\Urlmuskip=0mu plus 2mu\relax
\begin{itemize}[leftmargin=1.5em,itemsep=0.25em]
\item {\scriptsize\url{https://ai.google.dev/gemini-api/docs/changelog}}
\end{itemize}
\endgroup
\end{samepage}
\end{minipage}
\endgroup
\end{technicalnote}
\medskip
